\begin{lemma}
Assume \(\noderef{FiberAndDegreeEvent}\), the eventual-comparisons package
\(\noderef{ParameterHierarchyEventualComparisons}
(\varepsilon,\varepsilon_1,\varepsilon_2,n_0)\), \(n>n_0\), and
\[
|I|=\noderef{TwoBiteNaturalIndependenceScale}(1+\varepsilon,n).
\]
Then
\[
|\noderef{TwoBiteLargeHugeVertices}(\omega,\varepsilon,I)|<n^{1/4}.
\]
\end{lemma}
\begin{proof}
Fix \(n,m,p,\omega,I,\varepsilon,\varepsilon_1,\varepsilon_2,n_0\) satisfying
the hypotheses.  Put
\[
K=|I|=\noderef{TwoBiteNaturalIndependenceScale}(1+\varepsilon,n),\qquad
r=n^{1/4},\qquad
t_2=\noderef{TwoBiteLargeCutoff}(\varepsilon,n)=n^{1/4+\varepsilon},
\]
and set
\[
C_n=100(\log n)^3.
\]
Let
\[
A=\noderef{TwoBiteLargeHugeVertices}(\omega,\varepsilon,I)=L_I\cup H_I.
\]
Write \(a=|A|\).
Since \(n_0<n\) and \(n_0\) is a natural number, \(1\le n\), hence
\[
r=n^{1/4}\ge1>0.
\]
If \(A=\emptyset\), then \(a=0<r\) and the conclusion is immediate.  Thus the
only remaining case to exclude is the case where \(A\) is nonempty, and the
contradiction argument below will in fact choose a nonempty subfamily of \(A\).

We first record the pointwise lower bound for every \(x\in A\).  Unfolding
\(\noderef{TwoBiteLargeHugeVertices}\), membership \(x\in A\) means that either
\[
\noderef{TwoBiteLargeClass}(\omega,\varepsilon,I,x)
\quad\text{or}\quad
\noderef{TwoBiteHugeClass}(\omega,I,x)
\]
holds.  In the large case, the definition of
\(\noderef{TwoBiteLargeClass}\), together with
\(\noderef{TwoBiteLargeCutoff}\), gives directly
\[
|X_x(I)|>t_2.
\]
In the huge case, the definition of \(\noderef{TwoBiteHugeClass}\) gives
\[
\noderef{TwoBiteHugeCutoff}(n)<|X_x(I)|.
\]
Unfolding \(\noderef{TwoBiteHugeCutoff}\), this cutoff is exactly
\[
t_1=\frac{\sqrt{n\log n}}{\log\log n}.
\]
The eventual-comparisons hypothesis, applied to the present \(n>n_0\), contains
the explicit cutoff comparison
\[
\noderef{TwoBiteLargeCutoff}(\varepsilon,n)
<\noderef{TwoBiteHugeCutoff}(n).
\]
Since \(t_2=\noderef{TwoBiteLargeCutoff}(\varepsilon,n)\), the huge case also
implies \(|X_x(I)|>t_2\).  Hence, for all \(x\in A\),
\[
|X_x(I)|>t_2. \tag{1}
\]

Next, \(\noderef{TwoBiteX}\) says
\[
X_x(I)=I\cap N_x.
\]
Therefore \(X_x(I)\subseteq I\) for every \(x\).  If \(x,y\in A\) are distinct,
then
\[
X_x(I)\cap X_y(I)\subseteq N_x\cap N_y.
\]
We now match this containment to the three lifted-neighborhood intersection
clauses of \(\noderef{FiberAndDegreeEvent}\).  Write base vertices as elements
of \(V_R\cup V_B\).  If \(x=r_i\in V_R\) and \(y=r_j\in V_R\), then
\(x\ne y\) means \(i\ne j\), and the red-red lifted-neighborhood clause gives
\[
|N_{r_i}\cap N_{r_j}|\le100(\log n)^3=C_n.
\]
If \(x=b_i\in V_B\) and \(y=b_j\in V_B\), then \(i\ne j\), and the blue-blue
lifted-neighborhood clause gives
\[
|N_{b_i}\cap N_{b_j}|\le100(\log n)^3=C_n.
\]
If \(x=r_i\in V_R\) and \(y=b_j\in V_B\), or conversely, the mixed red-blue
lifted-neighborhood clause gives
\[
|N_{r_i}\cap N_{b_j}|\le100(\log n)^3=C_n.
\]
In each case the containment above gives
\[
|X_x(I)\cap X_y(I)|\le C_n \tag{2}
\]
for every distinct \(x,y\in A\).

We now spell out the finite inclusion-exclusion estimate.  Let
\(B\subseteq A\) be any finite subfamily, enumerate it as
\[
B=\{x_1,\ldots,x_b\},
\]
where \(b=|B|\), and put \(Y_i=X_{x_i}(I)\).  Since each \(Y_i\subseteq I\),
\[
K=|I|\ge\left|\bigcup_{i=1}^bY_i\right|.
\]
For these \(b\) sets one has
\[
\left|\bigcup_{i=1}^bY_i\right|
\ge
\sum_{i=1}^b |Y_i|
-
\sum_{1\le i<j\le b}|Y_i\cap Y_j|. \tag{3}
\]
Indeed, (3) follows by induction on \(b\).  The step from \(b\) to \(b+1\) uses
\[
|U\cup Y_{b+1}|=|U|+|Y_{b+1}|-|U\cap Y_{b+1}|,
\qquad U=\bigcup_{i=1}^bY_i,
\]
and the containment
\[
U\cap Y_{b+1}
\subseteq
\bigcup_{i=1}^b(Y_i\cap Y_{b+1}),
\]
so
\[
|U\cap Y_{b+1}|
\le
\sum_{i=1}^b |Y_i\cap Y_{b+1}|.
\]
Combining (1), (2), and (3) gives the following estimate for every finite
subfamily \(B\subseteq A\):
\[
K\ge
b\,t_2-\binom{b}{2}_{\mathbb R}C_n.
\]
If \(B\ne\emptyset\), then \(b>0\), and the first sum is a nonempty sum of
strict inequalities \(|Y_i|>t_2\).  In that case the estimate is strict:
\[
K>
b\,t_2-\binom{b}{2}_{\mathbb R}C_n. \tag{4}
\]
Here is the real-valued chain explicitly.  First cast all finite cardinalities
in (3) to \(\mathbb R\).  Since \(K=|I|\) as natural numbers and
\(\bigcup_iY_i\subseteq I\), the left side gives
\[
K\ge \left|\bigcup_{i=1}^bY_i\right|.
\]
If \(b>0\), the lower bound (1) gives, term by term after coercion to real
numbers,
\[
\sum_{i=1}^b |Y_i|> b\,t_2.
\]
When \(b=0\), the corresponding statement is the equality \(0=0\), giving
only the weak version above.  In the nonempty case used below, the displayed
strict inequality is valid.
For the intersections, (2) gives
\[
\sum_{1\le i<j\le b}|Y_i\cap Y_j|
\le \binom{b}{2}_{\mathbb R}C_n,
\]
because there are exactly \(b(b-1)/2\) index pairs
\((i,j)\) with \(1\le i<j\le b\), and
\[
\binom{b}{2}_{\mathbb R}
=\noderef{RealChooseTwo}((b:\mathbb R))
=\frac{(b:\mathbb R)((b:\mathbb R)-1)}2
\]
is the real cast of that pair count.  Thus
\[
\left|\bigcup_iY_i\right|
>
b\,t_2-\binom{b}{2}_{\mathbb R}C_n,
\]
which proves (4).  After casting the integer pair count to \(\mathbb R\), the
factor is
\[
\binom{b}{2}_{\mathbb R}=\noderef{RealChooseTwo}((b:\mathbb R)).
\]
In particular, when \(A\ne\emptyset\), taking \(B=A\) and viewing
\(a=|A|\) as a real number yields the real-valued lower bound
\[
K>
a\,t_2-\noderef{RealChooseTwo}(a)C_n. \tag{5}
\]
This is the promised conversion from the set-theoretic estimate to the
real-valued expression with \(a=|A|\).

It remains to rule out \(a\ge r\).  Suppose for contradiction that
\(a\ge r\).  Let \(A_{\mathbb N}=|A|\), so \(a=(A_{\mathbb N}:\mathbb R)\).
Since \(r\ge0\), let \(b=\lceil r\rceil\) be the least natural number whose
real cast is at least \(r\).  The inequality \(r\le a=(A_{\mathbb N}:\mathbb R)\)
therefore implies \(b\le A_{\mathbb N}=|A|\).  Hence there is a
\(b\)-element subfamily
\(B\subseteq A\) with
\[
|B|=b=\lceil r\rceil
\]
by choosing any \(b\) elements of the finite set \(A\).  Since \(b\) is a
ceiling of the real number \(r\), we have
\[
r\le b\le r+1.
\]
Also \(r\ge1\), so \(b\ge1\) and this chosen subfamily is nonempty.  Therefore
the strict version of (4) applies to this \(B\).  Define
\[
f(u)=u\,t_2-\noderef{RealChooseTwo}(u)C_n.
\]
The value of \(f(b)\) is exactly the right side of (4) for this
\(b\)-element subfamily, while \(f(r)\) is the real expression that appears in
the hierarchy gap.  We compare \(f(b)\) and \(f(r)\) without any limiting
argument.  Since
\(\noderef{RealChooseTwo}(u)=u(u-1)/2\), for real numbers
\(r\le u\le v\le r+1\) we have the exact identity
\[
\begin{aligned}
f(v)-f(u)
&=(v-u)t_2-\frac{C_n}{2}\bigl(v(v-1)-u(u-1)\bigr)\\
&=(v-u)\left(t_2-\frac{C_n}{2}(u+v-1)\right).
\end{aligned}
\tag{6}
\]
Here \(v-u\ge0\), and from \(u,v\le r+1\) we get
\[
\frac{u+v-1}{2}\le r+\frac12.
\]
Also \(n_0<n\) implies \(1\le n\), so \(\log n\ge0\) and hence
\[
C_n=100(\log n)^3\ge0.
\]
The eventual-comparisons hypothesis contains exactly the comparison needed to
make the parenthesized factor in (6) nonnegative.  Substituting
\(r=n^{1/4}\) and \(C_n=100(\log n)^3\), that listed comparison is
\[
\noderef{TwoBiteLargeCutoff}(\varepsilon,n)
\ge
100\left(n^{1/4}+\frac12\right)(\log n)^3,
\]
and the left side is \(t_2\).  Thus for every \(r\le u\le v\le r+1\),
\[
t_2-\frac{C_n}{2}(u+v-1)\ge
t_2-C_n\left(r+\frac12\right)\ge0.
\]
This use of the previous inequality is purely algebraic: multiplying
\((u+v-1)/2\le r+1/2\) by the nonnegative quantity \(C_n\) gives
\[
\frac{C_n}{2}(u+v-1)\le C_n\left(r+\frac12\right),
\]
and subtracting these quantities from \(t_2\) reverses the direction in the
expected way:
\[
t_2-\frac{C_n}{2}(u+v-1)
\ge t_2-C_n\left(r+\frac12\right).
\]
The hierarchy comparison then identifies the final lower bound with \(0\).
The identity (6) then gives \(f(v)\ge f(u)\) on this whole interval.  Applying
this with the real numbers \(u=r\) and \(v=b=\lceil n^{1/4}\rceil\), using
\(r\le b\le r+1\), gives \(f(b)\ge f(r)\).  No monotonicity outside this
one-unit interval is used.  Therefore
\[
f(b)\ge f(r)
=n^{1/4}t_2
 -\noderef{RealChooseTwo}(n^{1/4})100(\log n)^3. \tag{7}
\]
From (4) applied to \(B\), \(K>f(b)\).  Combining this with (7) gives
\[
K>
n^{1/4}t_2
 -\noderef{RealChooseTwo}(n^{1/4})100(\log n)^3. \tag{8}
\]
But the eventual-comparisons hypothesis
\(\noderef{ParameterHierarchyEventualComparisons}
(\varepsilon,\varepsilon_1,\varepsilon_2,n_0)\), applied to the present
\(n>n_0\) and the present
\[
K=\noderef{TwoBiteNaturalIndependenceScale}(1+\varepsilon,n),
\]
contains exactly the strict large-huge gap
\[
K<
n^{1/4}t_2
 -\noderef{RealChooseTwo}(n^{1/4})100(\log n)^3. \tag{9}
\]
The right side in (8) is literally the same real number as the right side in
(9).  Thus (8) says \(K\) is strictly larger than that real number, while (9)
says \(K\) is strictly smaller than it, an impossibility by transitivity of
the strict order on \(\mathbb R\).  Therefore \(a<r\), i.e.
\[
|\noderef{TwoBiteLargeHugeVertices}(\omega,\varepsilon,I)|<n^{1/4},
\]
as required.
\end{proof}
